\section{Post-Training：RLHF 与 Constitutional AI}

Pre-training 提供广泛语言和世界知识，但不会自动产生符合产品意图的对话行为。Post-training 把 instruction、preference、principle 和 tool policy 注入模型。课堂从 RLHF 讲到 Constitutional AI/RLAIF，重点不是把 human feedback 完全替换，而是扩大监督规模并让行为目标更明确、可迭代。

\subsection{RLHF：把偏好训练成 proxy reward}

\term{reinforcement learning from human feedback}（RLHF）通常包含 supervised fine-tuning、human preference comparison、\term{reward model}（奖励模型）和 policy optimization。Reward model 学习预测人类更喜欢哪个回答，RL 再优化模型提高 predicted reward，并用 KL 等约束避免策略偏离过远。

\lecturefigure{09-rlhf-pipeline.png}{RLHF 把 demonstrations、preference、reward model 与 policy optimization 串成训练管线。}{InstructGPT；本地字幕 21:00--26:20；概念重绘。}

\begin{knowledgebox}{读图：每一层都有独立误差}
SFT 受 demonstration 覆盖限制；preference 受 annotator/rubric 影响；reward model 可能在 distribution shift 下被 exploit；RL policy 可能提高 proxy reward 却损害真实 helpfulness。最终 human evaluation、safety evaluation 和 product metrics 不能被 reward curve 替代。
\end{knowledgebox}

\begin{warningbox}{Reward hacking 是 proxy 问题，不只是算法 bug}
任何有限 rubric 都无法完整表达“有帮助、诚实、无害且符合上下文”。模型可能学会讨好 judge、过度拒绝或使用表面风格获得高分。应使用多维 evaluation、adversarial examples、holdout tasks 和 policy review，避免单一 reward 成为唯一真理。
\end{warningbox}

\subsection{Constitutional AI 与 RLAIF}

\term{Constitutional AI} 用一组原则指导模型 critique 和 revise 自己的回答，再生成 AI preference data；\term{reinforcement learning from AI feedback}（RLAIF）用 evaluator model 产生偏好或 reward signal。人类仍负责选择 principles、审计行为、处理冲突与决定 deployment boundary，AI feedback 负责降低每个样本都要人工标注的成本。

\lecturefigure{10-constitutional-ai.png}{Constitutional AI 用 principles、critique/revision、AI preferences 与 RLAIF 扩展监督。}{Anthropic Constitutional AI；本地字幕 26:00--27:40；概念重绘。}

\begin{knowledgebox}{读图：Base model capability 是前提}
模型必须能理解原则、发现问题并提出更好 revision，CAI 才能工作；evaluator model 也可能共享 target model 的 blind spot。Principle 冲突、文化语境、过度拒绝与 hidden reasoning 都需要 human audit。RLAIF 扩大监督吞吐，不等于消除 human governance。
\end{knowledgebox}

\teachervoice{Mann 把 Constitutional AI 描述为一个能力门槛后的突破：早期模型还不足以可靠 critique 自己，模型变强后才可能成为监督工具。这个例子说明 capability growth 也能创造新的 safety technique，而不是只增加风险。}

\subsection{本章小结}

Post-training 是多层 proxy 优化。RLHF 和 RLAIF 都需要清晰原则、独立 evaluation 与 human governance；它们改变行为，却不替代 base capability、system control 或 deployment safety。

